\documentclass[a4paper,11pt]{article}

\usepackage[T1]{fontenc}
\usepackage[utf8]{inputenc}
\usepackage[margin=1.7cm]{geometry}
\usepackage{xcolor}
\usepackage{array}
\usepackage{longtable}
\usepackage[colorlinks=true,linkcolor=blue,urlcolor=blue,citecolor=blue]{hyperref}

\title{Potential Journals for Neurorobotics-Related Papers\\
\large Revised Shortlist Based on the NES Classification Criteria}
\author{Mohammad Rahmani}
\date{October 2026}

\newcolumntype{L}[1]{>{\raggedright\arraybackslash}p{#1}}

\begin{document}
\maketitle

\noindent
\textbf{Paper:} \emph{Self-awareness in Multi-Robot Systems}. The priorities below
reflect the paper's emphasis on neurally inspired cognitive architecture,
hierarchical internal language derived from sensory data, self--other distinction,
novelty detection, and autonomous multi-robot navigation.

\medskip
\noindent
\textbf{NES criterion:} Under the supplied NES classification grid, a
\emph{Science-to-Science} journal article is classified as Category~I if the
journal was listed in Q1 in at least one Scimago subject category in either
of the preceding two years. Q2 or Q3 journals fall under Category~II.

\renewcommand{\arraystretch}{1.2}
\setlength{\tabcolsep}{5pt}
\small

\begin{longtable}{@{}L{0.045\textwidth}L{0.245\textwidth}L{0.070\textwidth}L{0.065\textwidth}L{0.460\textwidth}@{}}
\caption{Recommended journals, ranked by suitability for the current paper.}
\label{tab:neurorobotics_journals_nes}\\
\hline
\textbf{Rank} & \textbf{Journal} & \textbf{SJR} & \textbf{NES} & \textbf{Suitability for this paper}\\
\hline
\endfirsthead
\multicolumn{5}{l}{\emph{Table \thetable\ continued from previous page}}\\
\hline
\textbf{Rank} & \textbf{Journal} & \textbf{SJR} & \textbf{NES} & \textbf{Suitability for this paper}\\
\hline
\endhead
\hline
\multicolumn{5}{r}{\emph{Continued on next page}}\\
\endfoot
\hline
\endlastfoot

1 & \href{https://cis.ieee.org/publications/t-cognitive-and-developmental-systems}{IEEE Transactions on Cognitive and Developmental Systems}
& Q1 & I & \textbf{Very high.} Cognitive architectures, developmental learning, and brain-inspired robotic cognition.\\

2 & \href{https://www.keaipublishing.com/en/journals/cognitive-robotics/}{Cognitive Robotics}
& Q1 & I & \textbf{Very high.} Cognitive processes and learning in robots; especially suitable for sensory-to-language representation learning.\\

3 & \href{https://link.springer.com/journal/10514}{Autonomous Robots}
& Q1 & I & \textbf{High.} Robotic learning, perception, autonomy, and adaptive behavior.\\

4 & \href{https://www.sciencedirect.com/journal/robotics-and-autonomous-systems}{Robotics and Autonomous Systems}
& Q1 & I & \textbf{High.} Cognitive architecture demonstrated through autonomous navigation and multi-robot experiments.\\

5 & \href{https://journals.sagepub.com/home/adb}{Adaptive Behavior}
& Q1\textsuperscript{*} & I\textsuperscript{*} & \textbf{High.} Biologically inspired adaptive behavior, cognitive mechanisms, and regulation.\\

6 & \href{https://www.sciencedirect.com/journal/neural-networks}{Neural Networks}
& Q1 & I & Most suitable if the principal novelty is the neural learning or representation algorithm.\\

7 & \href{https://link.springer.com/journal/12369}{International Journal of Social Robotics}
& Q1 & I & Suitable if robot--robot communication and self--other distinction are central contributions.\\

8 & \href{https://direct.mit.edu/neco}{Neural Computation}
& Q1\textsuperscript{*} & I\textsuperscript{*} & Suitable if the primary contribution is the computational or theoretical neural model.\\

9 & \href{https://www.frontiersin.org/journals/neurorobotics}{Frontiers in Neurorobotics}
& Q2 & II & \textbf{Very high topical fit,} but lower NES classification under the stated Q2 rankings.\\

10 & \href{https://www.sciencedirect.com/journal/cognitive-systems-research}{Cognitive Systems Research}
& Q2 & II & Very suitable for cognitive architectures and theoretical models, but lower NES classification.\\

11 & \href{https://www.ieee-ras.org/publications/ra-l/}{IEEE Robotics and Automation Letters}
& Q1 & I & Suitable for a focused robotics contribution supported by strong empirical results.\\

12 & \href{https://iopscience.iop.org/journal/1748-3190}{Bioinspiration \& Biomimetics}
& Q2 (2025); Q1 (2024) & I & Suitable if the biological inspiration is more clearly formalized and empirically validated.\\

\end{longtable}

\normalsize
\section*{Additional ambitious targets}
\begin{itemize}
  \item \href{https://www.ieee-ras.org/publications/t-ro/}{IEEE Transactions on Robotics}
  (Q1; NES~I): Consider if the robotics novelty and experimental validation
  become substantially stronger.
  \item \href{https://www.science.org/journal/scirobotics}{Science Robotics}
  (Q1; NES~I): A highly ambitious option requiring an exceptional and
  broadly significant contribution with convincing experimental evidence.
\end{itemize}

\section*{Classification notes}
\begin{itemize}
  \item \textsuperscript{*} For \emph{Adaptive Behavior}, the noted Q1 subject category
  is \emph{Philosophy} rather than \emph{Artificial Intelligence}. For
  \emph{Neural Computation}, the noted Q1 category is
  \emph{Arts and Humanities (miscellaneous)}. The supplied NES rule refers
  to Q1 in \emph{any} Scimago subject category, but administrative
  confirmation is advisable.
  \item For \emph{Bioinspiration \& Biomimetics}, the preceding shortlist reported
  Q1 in 2024 and Q2 in 2025. Under the stated two-year NES criterion, the
  2024 Q1 ranking would qualify it for Category~I if that year lies within
  the applicable evaluation window.
  \item These SJR quartiles reproduce the preceding journal shortlist;
  check the applicable Scimago years and subject categories again before submission.
  \item The priority ranking describes \emph{topic fit for the current manuscript},
  not a universal ranking of journal quality or acceptance probability.
\end{itemize}

\end{document}
